% !PRINTED 324
\weberpage{324}

%<<CONT
when $n$ is an odd prime (by the law of reciprocity and its supplementary theorems [Bd. I, § 145, 9.]), and these formulae can be summarised in
\weq{10}{(\pm p, n) = \left(\dfrac{n}{p}\right).}
But by (7) this formula holds for \gesperrt{every arbitrary $n$ that is not divisible by $p$.}

In the same way one finds, for an odd $n$:
\weq{11}{\begin{aligned} (-4,n) &= (-1)^{\frac{n-1}{2}},\\ (8,n) &= (-1)^{\frac{n^2-1}{8}},\\ (-8,n) &= (-1)^{\frac{n-1}{2}+\frac{n^2-1}{8}}. \end{aligned}}
Further, in general, if $D$, $D'$ are two discriminants:
\weq{12}{(D,n)(D',n) = (DD',n).}
This formula holds in the first place when $DD'$ and $n$ are not relatively prime, because then both sides are $=0$. It holds further for $n=-1$ [by (3)], for $n=2$ [by (5) and the congruence $\dfrac{D^2-1}{8}+\dfrac{D'^2-1}{8} \equiv \dfrac{D^2D'^2-1}{8}$ (mod 2)] and for an odd prime $n$ [by (6)]. Hence by (7) it holds in general.

We can now summarise the law of reciprocity together with its supplementary theorems for discriminants in the following way:
\weq{13}{(D,D') = \pm(D',D),}
\gesperrt{in which the upper sign holds when at least one of the two discriminants $D$, $D'$ is positive, the lower when both are negative.}

In order to derive this formula from the usual form of the law of reciprocity, we assume first that the two formulae hold:
\[
(D,D') = \pm(D',D),
\]
\[
(D_1,D') = \pm(D',D_1).
\]
Then from (8) and (12) it follows that
\[
(DD_1,D') = \pm(D',DD_1),
\]
in which again the lower sign holds only when $D'$ and $DD_1$ are both negative. By a repeated application there results
\[
(DD_1,D'D_1') = \pm(D'D_1',DD_1)
\]
%CONT>>
